\chapter{Kinship in India}

\section{The Feminist Critique of the Family}

\begin{KeyConcept}
\begin{itemize}
\item Feminist studies of the family came late to sociology. They argued the family is \textbf{not egalitarian} but a system that binds women (through marriage), and they introduced \textbf{domestic violence, unpaid (household) labour, and the sexual division of labour} into the discipline. (Durkheim regarded the \textbf{sexual division of labour as the first form of division of labour} --- but feminists point out that it is a discriminatory division.)
\item \textbf{Margaret Benston}: women's \textbf{unpaid labour is beneficial to capitalism} --- if housework were paid, wages would rise and workers could confront capital; even a minimum wage for women would be a massive redistribution. The (male) breadwinner's obligation to his family \textbf{delays conflict} (he cannot withdraw from wage labour).
\item \textbf{Fran Ansley}: the family works as a \textbf{'safety valve'} --- it absorbs the worker's alienation so industrialism is not threatened. \textbf{David Cooper} ('The Death of the Family'): family socialisation is repressive --- it makes children \textbf{obedient and submissive} to authority (the labour force capitalism needs); so the family must be remodelled.
\item Feminists also showed the family is a \textbf{political institution} --- an inequality of power between men and women (husband dominating wife).
\end{itemize}
\end{KeyConcept}

\begin{QuickRevision}
\begin{itemize}
\item not egalitarian
\item domestic violence, unpaid (household) labour, and the sexual division of labour
\item sexual division of labour as the first form of division of labour
\item Margaret Benston
\item unpaid labour is beneficial to capitalism
\item delays conflict
\item Fran Ansley
\item 'safety valve'
\item David Cooper
\item obedient and submissive
\item political institution
\end{itemize}
\end{QuickRevision}

\section{Marital Breakdown}

\begin{KeyConcept}
\begin{itemize}
\item \textbf{Marital breakdown} takes three forms: \textbf{divorce} (legal termination), \textbf{separation} (living apart without divorce), and \textbf{empty-shell marriages} (living together without a relationship).
\item \textbf{Nicky Hart} ('When Marriage Ends') --- three factors: the \textbf{value attached to marriage}, the \textbf{degree of conflict between spouses}, and the \textbf{opportunities to escape} marriage.
\item On value: \textbf{Parsons and Ronald Fletcher} argue the value of marriage has \textbf{increased} (rising remarriage and higher expectations --- people leave a marriage that fails their expectations). On conflict: \textbf{William Goode} (emotional over-strain on the isolated nuclear family, whose functions were taken by schools/hospitals), \textbf{Edmund Leach} (economic and social bonds weakened), \textbf{Graham Allan \& Graham Crow} (independent incomes and the decline of family-owned businesses make divorce possible). On escape: \textbf{Colin Gibson} (secularisation made divorce acceptable), divorce legislation (the Special Marriage Act, Hindu Marriage Act), and the rise of \textbf{serial monogamy} (divorce $\rightarrow$ remarriage $\rightarrow$ divorce).
\item The class dimension: divorce is acceptable in the urban/middle class but not the lower class; children are affected more by \textbf{parental conflict} than by divorce itself (which is socially acceptable in the West).
\end{itemize}
\end{KeyConcept}

\begin{QuickRevision}
\begin{itemize}
\item Feminist critique: unpaid labour (Benston), safety valve (Ansley), repressive socialisation (Cooper); family = political/power institution.
\item Marital breakdown: divorce, separation, empty-shell; Nicky Hart's 3 factors; Parsons/Fletcher (value up), Goode (over-strain), Leach (weak bonds), Allan \& Crow (independent income), Gibson (secularisation), serial monogamy.
\end{itemize}
\end{QuickRevision}

\section{Intimacy: From Spiritual Love to Plastic Sexuality}

\begin{KeyConcept}
\begin{itemize}
\item \textbf{Anthony Giddens} on the changing idea of intimacy: (1) \textbf{spiritual love} --- a pure, unattainable love that cannot become marriage (Krishna and Radha; Laila-Majnu); (2) \textbf{romantic love} --- a product of the Enlightenment novel, love that \emph{can} be achieved within social boundaries (same caste/status), committed for life, placed inside marriage; (3) \textbf{plastic sexuality} --- based on \textbf{confluent love} (stay while it satisfies, otherwise leave), focused on physical beauty (Tinder/Bumble), expressed in \textbf{serial monogamy} and alternate forms of intimacy.
\item Feminists partly favour plastic sexuality (it gives women choice), but note that only \textbf{independent women} (economically and socially) can practise it openly.
\end{itemize}
\end{KeyConcept}

\begin{QuickRevision}
\begin{itemize}
\item Anthony Giddens
\item spiritual love
\item romantic love
\item plastic sexuality
\item confluent love
\item serial monogamy
\item independent women
\end{itemize}
\end{QuickRevision}

\section{Iravati Karve: The Four Kinship Patterns}

\begin{KeyConcept}
\begin{itemize}
\item \textbf{Iravati Karve} (an Indologist, student of G. S. Ghurye, of the Deccan College, and author of the first feminist reinterpretation of the Mahabharata) studied kinship through scriptures. Kinship is of three types --- by \textbf{blood (consanguineal)}, by \textbf{marriage/adoption (affinal)}, and \textbf{fictive} (now powerful as families are isolated from kin --- a move from kinship to friendship).
\item Kinship patterns depend on \textbf{four factors --- language, geography, caste and kinship terminology} --- giving \textbf{four regional patterns}:
\item \textbf{North Indian} (Punjabi, Sindhi, Haryanvi, Hindi, Bihari): \textbf{patrilineal}, \textbf{gotra exogamy}, \textbf{sapinda} (5 generations of the father, 3 of the mother), \textbf{territorial/village exogamy}, \textbf{caste endogamy}, and \textbf{kanyadan} (donation --- the giver's status is low and he accepts nothing in return, sending gifts = dowry). Women as \emph{daughter} have no importance; women as \emph{wife} are important.
\item \textbf{South Indian}: \textbf{cross-cousin marriage and uncle--niece marriage} (the uncle younger than the mother) are allowed, based on the \textbf{exchange of women}; gotra exogamy is both-sided. Exceptions: the \textbf{Todas} (polyandry --- a woman married to all brothers, where the biological father is unimportant; social fatherhood matters) and the \textbf{Nairs of Kerala} (matrilineal).
\item \textbf{Central Indian} (Gujarat, Rajasthan, Maharashtra, MP): mixed --- higher castes follow the North, lower castes the South (Gujarat's Kathiawar tribes cross-cousin, while Patidars/Brahmins do not --- D. F. Pocock's hypergamy); the \textbf{Rajputs} marry 'sor' (focusing on physical beauty rather than caste, hypergamous --- the Kshatriya category is fluid, per K. M. Panikkar); the \textbf{Jats} follow four-gotra exogamy; the \textbf{Marathas} allow maternal cross-cousin marriage but not paternal.
\item \textbf{Eastern Indian}: very mixed --- the \textbf{Nagas} practise \textbf{marriage by capture} (the groom challenges the bride's father); the \textbf{Bonda} tribe prefer \textbf{divorced/older women} (women hold a privileged position).
\end{itemize}
\end{KeyConcept}

\begin{QuickRevision}
\begin{itemize}
\item Giddens: spiritual $\rightarrow$ romantic $\rightarrow$ plastic sexuality (confluent love).
\item Karve: kinship by blood/adoption/fictive; 4 factors (language, geography, caste, terminology); 4 patterns (North = patrilineal/kanyadan/gotra-sapinda exogamy; South = cross-cousin/uncle-niece/exchange; Central = mixed/Rajput sor/Jat 4-gotra; Eastern = Nagas marriage-by-capture, Bonda older/divorced women).
\item Exceptions: Todas (polyandry), Nairs (matrilineal).
\end{itemize}
\end{QuickRevision}

\section{From Family to Household}

\begin{KeyConcept}
\begin{itemize}
\item \textbf{K. M. Kapadia} defined the joint family (kin members, three or more generations, common property, one roof). \textbf{I. P. Desai} (Mahua, Gujarat) added the \textbf{'degree of jointness'} (economic + emotional) --- so a \textbf{moderately joint family} accepts the patriarch's decisions even when living apart.
\item \textbf{A. M. Shah} argued the category 'family' \textbf{excludes more people than it includes} (migrants, new family forms) --- so we should study the \textbf{household}: a \textbf{residential and domestic unit of one or more persons living under the same roof and eating from a single kitchen}. He gives \textbf{four types --- simple, complex, parental, and compound households} (built on the 'parental family' = parents + unmarried children).
\item The \textbf{household} category (used by Patricia Uberoi, Pauline Kolenda and feminists) includes single-parent, gay, lesbian and live-in families --- and policy (e.g. MGNREGA) now targets households, not families.
\end{itemize}
\end{KeyConcept}

\begin{QuickRevision}
\begin{itemize}
\item Kapadia (joint family), I. P. Desai (degree of jointness), A. M. Shah (household --- 4 types; family excludes more than includes); household used by Uberoi/Kolenda and policy (MGNREGA).
\end{itemize}
\end{QuickRevision}
